\scanpage{53}
\mspage{9}
\uline{NB} Ayant (via un groupoïde équivalent à description purement combinatoire) ``réconstruit''
\deleted{\unsure{ou|du}} \deleted{\unsure{plutôt|purement}} réalisé le groupoïde
$S\mathcal TD_{\mathrm{top}}^\pm$, et le foncteur
\[
 S\mathcal TD_{\mathrm{top}}^\pm\longrightarrow(\mathrm{Gr.\ pond.\ rep})
\]
\deleted{\unsure{alors}} on aura \deleted{\unsure{par}} donc aussi \deleted{\unsure{une description}} \inserted{l'équivalent de}
\deleted{de} la donnée d'un $(X,\Sigma,R)$, muni \uline{en plus} d'une partie $R'\subset R$ --- \deleted{d'où} i.e. avec la construction d'un groupoïde \unsure{équivalent|au-dessus}
\diagram{\begin{tikzcd}[ampersand replacement=\&,row sep=large]
S_*\mathcal TD_{\mathrm{top}}^\pm\,
 \begin{gathered}(\text{à sous ens. marqué }R')\\
 \text{de sommets-}\\\text{repères}\end{gathered}\arrow[d]\\
S_*\mathcal TD_{\mathrm{top}}^\pm
\end{tikzcd}}
et \deleted{\unsure{d'où}} il nous faudra donc \inserted{décrire} un foncteur ``gommage de $R'$''
\diagram{\begin{tikzcd}[ampersand replacement=\&,column sep=huge]
S_*\mathcal TD_{\mathrm{top}}^\pm
 \left(\underset{\text{marqué}}{\overset{\text{ens}}{\deleted{\unsure{R_0|R}}R'}}\right)
 \arrow[r,"\text{gommage de }R'"]\arrow[r,shorten >=4em]
 \&S_*\mathcal TD_{\mathrm{top}}^\pm
\end{tikzcd}}
\deleted{Cela s'exprime, sur le}
donnant lieu : un diagramme comm. de foncteurs
\diagram{\begin{tikzcd}[ampersand replacement=\&,row sep=large,column sep=large]
S_*\mathcal TD_{\mathrm{top}}^\pm(\ldots R'\ldots)
 \arrow[rr,"\text{gommage de }R'"]
 \arrow[d,"\begin{gathered}\text{oubli du }\deleted{\unsure{R'}}\text{ marquage}\\R'\end{gathered}"]
 \&\&S_*\mathcal TD_{\mathrm{top}}^\pm\arrow[ddl]\\
S_*\mathcal TD\deleted{\unsure{\#}}_{\mathrm{top}}^\pm\arrow[dr]\&\&\\
\&(\mathrm{Gra}\deleted{\unsure{\mathrm{bi}}}\mathrm{pond})
 \arrow[r]\&
 \begin{gathered}\text{graphes avec}\\\text{pondération}\\
 \text{par }\gamma,\text{ mais}\\\text{non par }r,\end{gathered}
\end{tikzcd}}
alors qu'au niveau

\margin{\uline{NB} Cette opération a l'air d'un \uline{affaiblissement} de structure --- mais cet affaiblissement, \unsure{sur|pour} une composante $C$ donnée de $\Sigma^+$, lorsqu'on ``gomme'' \uline{tous} les repères de \deleted{\unsure{ce}} $C$ --- sinon c'est au contraire plutôt dans le sens d'un renforcement de structure, --- ou plutôt, des \unsure{marquages|groupoïdes} \unsure{qui} \unsure{du nombre de|du même type} $R'_C=R'\cap C$, dans le cas catégories isotopiques \deleted{\unsure{au dessus de}} \inserted{l'opération de gommage et celles} de surmarquage (\unsure{cf.} plus bas) \unsure{correspondantes} sont inverses \unsure{l'une de l'autre|l'un de l'autre}, sont des équivalences de groupoïdes.
Il faudrait séparer les deux cas, pour distinguer ces \uline{opérations de gommage essentielles} (qui changent le groupe des autom.) et les autres.}
\mspage{52}
